% ECONOMICS_PROBLEM: {"id": "H26-608", "title": "Digital tax administration and compliance", "jel_code": "H26", "source_id": "OP-0047"}
% !TeX program = xelatex
\documentclass[11pt,a4paper]{article}
\usepackage[margin=23mm]{geometry}
\usepackage{fontspec}
\setmainfont{Latin Modern Roman}
\usepackage{amsmath,amssymb,mathtools}
\usepackage{xcolor,enumitem,booktabs,longtable,array,xurl,hyperref}
\definecolor{muted}{HTML}{53636C}
\hypersetup{colorlinks=true,urlcolor=blue,linkcolor=blue}
\setlength{\parindent}{0pt}
\setlength{\parskip}{5pt}
\newcommand{\R}{\mathbb R}
\newcommand{\E}{\mathbb E}
\newcommand{\N}{\mathbb N}
\newcommand{\ind}{\mathbf 1}
\newcommand{\Var}{\operatorname{Var}}
\newcommand{\Cov}{\operatorname{Cov}}
\newcommand{\supp}{\operatorname{supp}}
\DeclareMathOperator*{\argmin}{arg\,min}
\DeclareMathOperator*{\argmax}{arg\,max}
\newcommand{\entrymeta}[3]{{\sffamily\small\color{muted}\textbf{\texttt{#1}}\quad|\quad #2\par #3\par}}
\newcommand{\Problemref}[1]{\texttt{#1}}
\newcommand{\JELref}[2]{\texttt{#2} (JEL \texttt{#1})}
\begin{document}
\section*{Digital tax administration and compliance}
\textbf{Problem ID:} \texttt{H26-608}\quad \textbf{JEL:} \texttt{H26}\par
% BEGIN ECONOMICS STATEMENT
\label{problem:OP-0047}
\label{audited:ECON-089}
\label{atlas:prob:P7}
\noindent\textbf{Original atlas identifier:} P7; entry 47. \textbf{Field:} Public finance and taxation. \textbf{Original tractability label:} Easy.

\noindent\textbf{Classification:} Parameterized theoretical/empirical research agenda; not a certified unsolved theorem.

\subsection*{Definitions and assumptions}
Let $d=(d_1,\ldots,d_K)$ represent digital-administration features such as e-invoicing, pre-filling, reporting latency, and risk-score auditing. A feasible policy includes an implementation schedule and audit capacity. Firm $i$ belongs to a transaction network $\mathcal G$ and has treatment exposure $h_i(d,\mathcal G)$; outcomes can depend on neighbors' assignments. Define net revenue $R(d)$ after refunds and corrected assessments, administrative resource cost $C_A(d)$, private compliance resource cost $C_B(d)$, and expected loss $C_E(d)$ from erroneous enforcement and data breaches. The last term requires an explicitly chosen monetary valuation and event-probability model. Let $EV_i(d)$ capture private welfare changes excluding the three separately counted costs; fix nonnegative social weights $w_i$ and fiscal-value parameter $\lambda$.

\subsection*{Formal research question}
\emph{Editorial formulation: the mathematical target below is not a theorem quoted from the references.}

Identify the causal policy value
\[
 W(d)=\mathbb E[w_iEV_i(d)]+\lambda R(d)-C_A(d)-C_B(d)-C_E(d)
\]
over feature combinations satisfying budget, capacity, correction/appeal, and legal-data-use constraints. Costs and equivalent variation are expressed in a common monetary welfare unit; count each resource or loss once. Determine which randomized feature assignments identify direct effects, network spillovers, and interactions rather than merely the effect of a bundled rollout. Characterize $\varepsilon$-optimal policies and bounds for unobserved compliance burdens or rare error losses. If feature assignment is staggered rather than randomized, state counterfactual-trend, anticipation, and interference assumptions before attributing the revenue change to digitization.

\subsection*{Interpretation and scope}
Higher assessed liability is not necessarily higher correct tax collection: include appeal reversals, payment delays, and displaced reporting in the revenue measure. “Lowest social cost” requires a service or revenue target, since doing nothing may minimize expenditure. Digital reporting can save firms time while increasing oversight; both convenience and enforcement mechanisms must be separately measured. Rare breaches cannot be estimated precisely from a short experiment, so the policy comparison needs sensitivity ranges rather than a fictitious zero cost. Supply-chain exposures make ordinary no-interference assumptions questionable; network or saturation randomization offers a more coherent design. The atlas's Easy label is retained as an original judgment, not a demonstrated tractability result for this joint welfare problem.

\subsection*{Source audit and status}
[S1] and [S2] provide information/reporting evidence; [S3] is an OECD design framework rather than a causal evaluation. Its publication year is 2020 despite later webpage migrations. The expanded cost-inclusive feature-selection problem is editorial; the cited sources do not certify a unique current optimum.

\subsection*{References}
\begin{enumerate}[label={[S\arabic*]},leftmargin=*,itemsep=0.6em]
\item Dina Pomeranz (2015). No Taxation without Information: Deterrence and Self-Enforcement in the Value Added Tax. American Economic Review 105(8), 2539--2569. DOI: 10.1257/aer.20130393.
\url{https://www.aeaweb.org/articles?id=10.1257/aer.20130393}
\newline\emph{Locator and support limits:} Abstract and experiments: Chilean VAT information and supply-chain enforcement; no comprehensive privacy-cost estimate.
\item Joana Naritomi (2019). Consumers as Tax Auditors. American Economic Review 109(9), 3031--3072. DOI: 10.1257/aer.20160658.
\url{https://pubs.aeaweb.org/doi/10.1257/aer.20160658}
\newline\emph{Locator and support limits:} Abstract: consumer reporting incentives in Brazil; not direct evidence for every digital feature.
\item OECD (2020). Tax Administration 3.0: The Digital Transformation of Tax Administration. Paris: OECD Publishing. DOI: 10.1787/ca274cc5-en.
\url{https://www.oecd.org/en/publications/tax-administration-3-0-the-digital-transformation-of-tax-administration_ca274cc5-en.html}
\newline\emph{Locator and support limits:} Report: institutional digital-transformation framework; descriptive/design guidance, not a causal cost-benefit estimate.
\end{enumerate}

\subsection*{Integrated refinement: ECON-089}
{\sffamily\small\color{muted}Digital tax data quality and revenue administration\par}
This refinement adopts the additional source conventions
(page~\pageref{audited:conventions}), subject to its local restrictions.
\textbf{Record coverage and quality versus real activity.}
Under digital reform $z$, let true activity be $Y_i(z)$, its record $Z_i(z)$,
and usable linkage indicator $M_i(z)$. Use an independently specified audit
selection/error model. Alongside net revenue
$\tau_R=\mathbb E[R(1)-R(0)]-\Delta C_{\rm admin}$, identify coverage changes and
a fixed-cohort quality contrast such as
\[
\tau_Q=\mathbb E[(Z_i(1)-Y_i(1))^2-(Z_i(0)-Y_i(0))^2].
\]
Require both records by design or declare a missing-record loss; do not silently
condition on post-reform linkage. Use common measurement units and finite
moments. Separate improved recording from real-activity, selection and
enforcement changes, including transaction-network spillovers and incorrect
assessments. An audit is a validation measurement with its own errors rather
than an assumed perfect truth label.
\subsection*{Supplied source audit for ECON-089}
Local [S1], [S2], etc. in this audit and the references immediately below
belong to ECON-089, independently of the original entry's reference numbering.
[S1] and [S2] support digitization, registry quality and revenue administration. This fixed-cohort/policy-dependent truth formulation is editorial and corrects the original estimand.
\subsection*{References for integrated refinement ECON-089}
{\footnotesize\raggedright
\par\noindent\textbf{[S1]} Anders Jensen, Anne Brockmeyer, and Lucie Gadenne (2024). \emph{Taxation and Development}. VoxDevLit 12(1), September 2024.\par \textit{Verified locator / source role:} Primary publisher page checked for review identity, authors, version and background. The specific published agenda is corroborated in [S2]; the exact formal target is editorial.\par \url{https://voxdev.org/voxdevlit/taxation}\par\smallskip
\par\noindent\textbf{[S2]} Oliver Hanney / VoxDev (living compilation). \emph{Evidence gaps: Key unanswered questions in development economics}. Accessed 8 October 2026. The page currently covers 20 topics; the ``16-topics URL is a historical slug.\par \textit{Verified locator / source role:} Topic heading: Taxation and Development. Supports the broad research agenda; this is not a verbatim quotation or an attribution of the displayed mathematical target.\par \url{https://voxdev.org/topic/evidence-gaps-key-unanswered-questions-across-16-topics-development-economics}\par\smallskip
}

\subsection*{Common conventions: Source mathematical conventions}

Unless an entry states otherwise, agent and firm sets are finite. State and action spaces are measurable, strategies and outcome maps are measurable, probability laws are specified, and expectations exist for the targets under discussion. Infinite-horizon discount factors lie in $(0,1)$ and discounted objectives are finite. Norms, tolerances and complexity input sizes must be specified for computational comparisons. Constants or primitives that appear locally are inputs, not empirical facts asserted by this catalogue.

\textbf{Identification.} A statistical model $m\in\mathcal M$ implies an observable law $P_m$ and target $\theta(m)$. At law $P$, its identified set is
\[
\Theta(P;\mathcal M)=\{\theta(m):m\in\mathcal M,\ P_m=P\}.
\]
Point identification holds when this set is a singleton, or equivalently when $P_{m_1}=P_{m_2}$ implies $\theta(m_1)=\theta(m_2)$ for all compatible models. Sharp bounds mean bounds attained, or approached when the boundary is not attained, by compatible models. Writing this set is a definition, not a solution: the substantive task is to establish informative restrictions and derive or estimate the set. An unrestricted-confounding model may give only trivial bounds.

\textbf{Inference.} Identification, estimability and valid uncertainty quantification are separate questions. Fix $\alpha\in(0,1)$. A confidence procedure $C_n$ over a declared law class $\mathcal P$ has asymptotic uniform coverage at level $1-\alpha$ when
\[
\liminf_{n\to\infty}\inf_{P\in\mathcal P}\Pr_P\{\theta(P)\in C_n\}\geq1-\alpha.
\]
For set-identified targets the coverage event must be defined separately, for example coverage of the full identified set. If an entry uses identification-indexed models instead of uniquely indexed laws, the same coverage convention is applied over those models. Pointwise limits do not establish uniform validity, and asymptotic validity does not guarantee finite-sample precision. Sample sizes, dependence structures and weak-identification sequences are part of each application.

\textbf{Counterfactuals.} $Y_i(a)$ is a potential outcome under a specified intervention, including its timing, implementation and versions. It is not $Y_i$ conditional on voluntarily choosing $a$. A full intervention vector permits interference; shorthand scalar treatment notation does not silently impose no interference. Assignment, selection, attrition, anticipation and measurement must be specified. A claim about an instrument's local effect is not automatically a population policy effect.

\textbf{Economic equilibrium.} A counterfactual model includes best responses, constraints, feasibility, market clearing, state transitions and expectations consistent with the stated information. When several equilibria exist, either a declared selector is an input or the target is set-valued. Comparisons across policy regimes must use the stated selection convention. Reduced-form relationships are not presumed invariant after an equilibrium-changing intervention.

\textbf{Optimization and welfare.} Optimization is over the stated feasible set. If attainment is not established, $\sup$ or $\inf$ and $\varepsilon$-optimal choices replace an asserted maximizer or minimizer. For example, with value $V=\sup_{a\in\mathcal A}W(a)<\infty$, an $\varepsilon$-optimal set is $\{a\in\mathcal A:W(a)\geq V-\varepsilon\}$ for $\varepsilon>0$. Benefits, costs, risk and health or environmental outcomes can be aggregated only using declared units and conversion weights. Interpersonal utility weights, the treatment of fiscal transfers and foreign profits, discounting, and the behavioral welfare criterion are explicit inputs. A comparative-static derivative additionally requires existence and appropriate differentiability of the selected counterfactual.

\textbf{Sources.} Local labels [S1], [S2], and so forth restart in every question. Locators identify the relevant abstract, model, section or result at the level actually checked; a bibliographic or abstract-level verification is not represented as inspection of an entire proof. Working papers are distinguished from later publications and changing versions. Source summaries are paraphrased. All URL checks and editorial formulations refer to the audit date stated above.

\subsection*{Common conventions: Additional-source status and mathematical conventions}

\label{audited:conventions}

Of the 100 supplied ECON entries, 65 distinct problems appear here; 32 close
matches refine earlier entries, and three duplicate questions map to existing
entries. Appendix~\ref{app:audited-crosswalk} lists every destination.
The attachment's P/R/S/U distinctions and source audits are retained. Its
precise mathematical formalizations are editorial unless the individual audit
says otherwise. Candidate subsidy targets ECON-004--006 retain their U flags;
the resolved approval-core question ECON-007 updates OP-0202 above.
The following supplied conventions govern both this part and the attributed
ECON refinements elsewhere in the catalogue, subject to local restrictions.

\textbf{Parameterized questions.} A problem family lists inputs that must be fixed before an instance is evaluated: population, horizon, policies, information, data law, model class and welfare weights. These are required choices, not quantities the citations are assumed to specify. Undefined applications should not be treated as identified estimands.

\textbf{Indices and integrability.} Sets of agents, firms and regions are finite unless a population measure is explicitly used. \(i,j\) index units, \(r\) regions, \(t\) dates and \(h\) horizons. \(\mathbb E\) and \(\Pr\) refer to the declared population and law; \(\mathbf1\{A\}\) is an indicator. Every displayed expectation and discounted sum needs existence in the stated domain. Infinite horizons use a discount in \((0,1)\) and a convergence condition; finite horizons may use discount one.

\textbf{Optimization and existence.} A displayed \(\max\), \(\min\), or \(\arg\max\) is conditional on attainment. For finite feasible menus this is immediate; general spaces need continuity and compactness/coercivity or another existence argument. Without attainment, the target is the supremum/infimum and arbitrarily near-optimal feasible choices. \(\inf\varnothing=+\infty\). Empty feasibility and an unidentified target are different issues.

\textbf{Potential outcomes.} \(Y_i(z)\) is the outcome under a specified intervention, not the observed conditional mean among people choosing \(z\). In binary comparisons \(\Delta Y_i=Y_i(1)-Y_i(0)\). Specify assignment, treatment versions, compliance, information and observation horizon. Interference is allowed under a full policy vector; compressed exposure notation needs a justified mapping. No interference, consistency and overlap are substantive assumptions, not automatic consequences of notation.

\textbf{Populations and observation.} Fix baseline eligibility and population weights. Conditioning on post-policy employment, survival, relocation or program uptake can change the population being compared. Death is not ordinary loss to follow-up; outcomes truncated by death need a composite convention, joint survival target, or explicitly defined survivor stratum. Fertility-changing interventions need a population functional rather than an unexplained fixed descendant cohort. Measurement and missingness models must be supplied.

\textbf{Identification.} A structural model \(m\in\mathcal M\) implies observable law \(P_m\) and target \(\theta(m)\). Identification means
\[
 P_{m_1}=P_{m_2}\ \Longrightarrow\ \theta(m_1)=\theta(m_2)
 \quad\text{for all }m_1,m_2\in\mathcal M .
\]
Without identification, the target is
\[
 \Theta(P;\mathcal M)=\{\theta(m):m\in\mathcal M,\ P_m=P\}.
\]
Writing \(\theta(P)\) before establishing identification can hide incompatible latent models. Confidence sets additionally need a sampling law, confidence level, asymptotic sequence and uniformity class. Point identification, coverage and informative precision are separate properties.

\textbf{Mediation.} Total effects of randomized assignment are distinct from cross-world natural indirect effects. Belief/effort-channel effects require a meaningful mediator intervention and additional measurement, exclusion or mediation assumptions. Randomization of the initial policy alone does not supply those assumptions. Report bounds or interventional alternatives when the stronger target is unsupported.

\textbf{Equilibrium.} A counterfactual \(W(z)\), \(p(z)\), or \(\mu_t^z\) requires individual optimization, feasibility, government/resource budgets, market clearing and state-transition laws. State equilibrium existence and selection, or report ranges over compatible equilibria. Initial-state integration includes subsequent endogenous transitions; current-state integration must use the policy-dependent distribution. Reduced-form invariance after policy is not assumed.

\textbf{Welfare accounts.} Scores, utility, health, dollars and ecological quantities require explicit conversion before addition. Fees, taxes, subsidies, wages and extortion payments are transfers between parties; distributional weights can make them welfare relevant, but their face value is not automatically a resource loss. State ownership, geographic boundaries, foreign-income weights and fiscal finance. Do not add profits to household welfare already including dividends, or count land capitalization and the underlying benefit twice. A benefit-cost expression must distinguish gross benefits from net welfare and subtract every real cost exactly once.

\textbf{Derivatives and risk.} Log derivatives require positive arguments; local responses require admissible perturbations and specified equilibrium branches. Differentiation under expectations needs regularity/domination, and boundaries may allow only one-sided derivatives. Risk uses a specified law; ambiguity uses a declared nonempty law set on a common history space. Adaptive policies must be nonanticipative. A robust criterion is an editorial choice unless attributed explicitly.

\textbf{Scope overlaps.} Allocation, collective choice, contracts, household bargaining and coordinated policy overlap economics with optimization and game theory. They are retained because they appeared in the original catalogue; no claim of a strictly game-theory-free selection is made.
% END ECONOMICS STATEMENT
\end{document}
